
\chapter*{Expos\'e \texorpdfstring{VI\textsubscript{B}}{VI_B}\\
Generalities on Group Schemes}
\phantomsection
\label{sga3:VIB:title}
\addcontentsline{toc}{chapter}{Expos\'e VI-B. Generalities on Group Schemes}
\markboth{EXPOS\'E VI\textsubscript{B}. GENERALITIES ON GROUP SCHEMES}%
  {EXPOS\'E VI\textsubscript{B}. GENERALITIES ON GROUP SCHEMES}

\begin{center}
  \textit{by J.-E. Bertin}
\end{center}

This Expos\'e, which does not correspond to any oral lecture in the
seminar, is intended to bring together a number of commonly used
technical results concerning group schemes.%
\begingroup
\renewcommand{\thefootnote}{\fnsymbol{footnote}}%
\footnote[1]{This Expos\'e was quite substantially revised after its
mimeographed edition; in particular, \hyperref[sec:sga3-VIB-10]{Sections 10} and 11 were completely
rewritten.}%
\endgroup

\section*{1. Morphisms of groups locally of finite type over a field}
\phantomsection
\label{sec:sga3-VIB-1}
\addcontentsline{toc}{section}{1. Morphisms of groups locally of finite
  type over a field}

\medskip
\phantomsection\label{par:sga3-VIB-1.1}%
\noindent\textbf{1.1.---}
Let \(A\) be an Artinian local ring, let \(G\) and \(H\) be two
\(A\)-groups,%
and let \(u:G\to H\) be a morphism of \(A\)-groups. Then \(u\)
induces a group morphism
\[
  u(A):G(A)\longrightarrow H(A).
\]%
Since \(H(A)\) acts on \(H\) by right translation, restriction along
\(u(A)\) defines an action of \(G(A)\) on \(H\). This action is
compatible with \(u\) and with the action of \(G(A)\) on \(G\) by
right translation. Since \(G(A)\) acts transitively on the strictly
rational points of \(G\) (see Expos\'e~VI\(_A\), \SGARefLink{sga3:VIA:0.4}{0.4} for the
definition), these points ``all behave in the same way with respect
to \(u\).'' The following properties spring from this observation.

\medskip
\phantomsection\label{prop:sga3-VIB-1.2}%
\noindent\textbf{Proposition 1.2.---}
\emph{Let \(u:G\to H\) be a quasi-compact morphism between
\(A\)-groups locally of finite type over \(A\). Then the set
\(u(G)\) is closed in \(H\), and}
\[
  \dim G=\dim u(G)+\dim\ker u.
\]%

Since \(u\) commutes with the inversion morphisms of \(G\) and \(H\),
the image \(u(G)\) is invariant under inversion on \(H\); the same is
therefore true of its closure \(\overline{u(G)}\) in \(H\). On the
other hand, let \(L\) be the set of points of \(H\times_A H\) whose
two projections belong to \(u(G)\). It is clear that \(L\) is the
image of
\[
  u\times_Au:G\times_A G\longrightarrow H\times_A H.
\]
The multiplication morphism of \(H\) consequently sends \(L\) into
\(u(G)\); in other words,
\[
  u(G)\cdot u(G)=u(G).
\]
\hyperref[lem:sga3-VIB-1.2.1]{Lemma~1.2.1} below shows that the closure \(\overline L\) of \(L\)
in \(H\times_AH\) is the set of points whose two projections belong
to \(\overline{u(G)}\). Hence
\[
  \overline{u(G)}\cdot\overline{u(G)}
    =\overline{u(G)}.
\]
Thus the reduced subscheme of \(H\) whose underlying space is
\(\overline{u(G)}\) is naturally equipped with a group structure in
the category \((\mathbf{Sch}_{/k})_{\mathrm{red}}\), where \(k\) is
the residue field of \(A\) (cf.\ Expos\'e~VI\(_A\), \SGARefLink{sga3:VIA:0.2}{0.2}).

Let us prove the first assertion of \hyperref[prop:sga3-VIB-1.2]{1.2}. After replacing \(A\) by
an algebraic closure of its residue field \(k\), we may suppose that
\(A\) is an algebraically closed field \(k\)
(cf.\ EGA~IV\(_2\), 2.3.12). Replacing \(u\) by
\[
  u_{\mathrm{red}}:G_{\mathrm{red}}\longrightarrow H_{\mathrm{red}},
\]
we may suppose that \(G\) and \(H\) are reduced. In this case, as we
have just seen, \(u(G)\) is the underlying space of a reduced group
subscheme of \(H\);%
\begingroup
\renewcommand{\thefootnote}{\fnsymbol{footnote}}%
\footnote[1]{\label{note:sga3-VIB-1.2-image-ambient}Translator's
note: the French source prints ``a reduced group subscheme of \(G\)''
at this point. Since \(u(G)\subset H\), and the preceding paragraph
constructs the reduced subscheme inside \(H\), the ambient group is
\(H\).}%
\endgroup{}
we may therefore suppose that \(u\) is dominant. Then \(G(k)\) acts
transitively on the set of connected components of \(H\), and it is
enough to prove that \(u(G)\cap H^0\) is closed. We are thus reduced
to the case in which \(H\) is connected, hence irreducible and of
finite type (Expos\'e~VI\(_A\), \SGARefLink{cor:sga3-VIB-2.4}{2.4}). The morphism \(u\) is then of
finite type, since it is quasi-compact and locally of finite type.
As \(H\) is Noetherian, \(u(G)\) is constructible
(EGA~IV\(_1\), 1.8.5), and therefore contains an open subset \(V\)
of \(H\) (EGA~0\(_{\mathrm{III}}\), 9.2.2). By
Expos\'e~VI\(_A\), 0.5, we then have
\[
  H=V\cdot V\subset u(G)\cdot u(G)=u(G).
\]

Let us prove the second assertion. Recall first that the functor
\(\ker u\) (cf.\ \SGARefLink{sga3:i:definition:2-3-6-1}{Expos\'e~I, 2.3.6.1}) is represented by
\(u^{-1}(e)\), where \(e\) denotes the identity element of \(H\).
When \(u\) is locally of finite type, \(\ker u\) is therefore
locally of finite type over \(A\). As above, this time using
EGA~IV\(_2\), 4.1.4, we reduce to the case in which \(A\) is an
algebraically closed field \(k\).

We may further suppose that \(G\) and \(H\) are irreducible and of
finite type and that \(u\) is dominant. Indeed, since \(k\) is
algebraically closed, the connected components of \(G\), on whose
set \(G(k)\) acts transitively, clearly all have the same dimension;
and if \(u^0\) denotes the restriction of \(u\) to \(G^0\), then
\[
  (\ker u)^0\subset\ker u^0,
  \qquad
  \dim\ker u^0=\dim\ker u.
\]
Likewise, \(u\) is then of finite type over \(k\). If \(\eta\) is
the generic point of \(H\), then
\[
  \dim u^{-1}(\eta)=\dim G-\dim H
\]
(EGA~IV\(_3\), 10.6.1(ii)). By EGA~IV\(_3\), 9.2.3 and 9.2.6, the
set of \(y\in H\) such that
\[
  \dim u^{-1}(y)=\dim u^{-1}(\eta)
\]
contains a nonempty open subset \(V\). Since \(u\) is dominant,
\(U=u^{-1}(V)\) is a nonempty open subset of \(G\); it contains a
closed point \(x\), because \(G\) is a Jacobson scheme
(EGA~IV\(_3\), 10.4.7). Right translation \(r_x\) is an isomorphism
from \(\ker u\) onto \(u^{-1}(u(x))\), and hence
\[
  \dim\ker u
    =\dim u^{-1}(u(x))
    =\dim u^{-1}(\eta)
    =\dim G-\dim H.
\]

\medskip
\phantomsection\label{lem:sga3-VIB-1.2.1}%
\noindent\textbf{Lemma 1.2.1.---}
\emph{Let \(f:X'\to X\) and \(g:Y'\to Y\) be two quasi-compact
dominant morphisms of schemes over an Artinian local ring \(A\).
Then}
\[
  f\times_Ag:X'\times_A Y'\longrightarrow X\times_A Y
\]
\emph{is dominant (and quasi-compact).}

Indeed,
\[
  f\times_Ag
    =
  (\operatorname{id}_X\times_Ag)
    \circ
  (f\times_A\operatorname{id}_{Y'}).
\]
It is therefore enough to prove that
\(f\times_A\operatorname{id}_{Y'}\) and
\(\operatorname{id}_X\times_Ag\) are dominant. For this purpose we
may replace \(A\) by its residue field \(k\). Then \(X\) and \(Y'\)
are flat over \(A=k\), and
\(f\times_A\operatorname{id}_{Y'}\)
(respectively, \(\operatorname{id}_X\times_Ag\)) is obtained from
\(f\) (respectively, \(g\)) by the flat base change \(Y'\to A\)
(respectively, \(X\to A\)). It is therefore dominant (and
quasi-compact) by EGA~IV\(_2\), 2.3.7.
